\section{An elliptic boundary example}\label{sec:elliptic}

We prove Theorem~\ref{thm:elliptic} by rational identities. No
elliptic-curve group law is needed. Work first in the function field
\[
 F=\Q(x,y),\qquad y^2=x^3-2,
\]
where $x$ is transcendental. Define the rational map
\begin{equation}\label{eq:elliptic-map}
 \begin{aligned}
 x'&=\frac{\cA(x,y)}{(x-3)^2},
 &\cA(x,y)&=3x^2+9x-4-10y,\\
 y'&=\frac{\cB(x,y)}{(x-3)^3},
 &\cB(x,y)&=5x^3+45x^2-40-(27x+19)y.
 \end{aligned}
\end{equation}
Write $(x_i,y_i)$ for its iterates, starting with $(x_0,y_0)=(x,y)$.
The circle parameters will be $x_0,\ldots,x_6$.

\subsection{A finite rational test and real embeddings}
The polynomial identity
\begin{equation}\label{eq:curve-identity}
 \cB(x,y)^2-\cA(x,y)^3+2(x-3)^6
 =(y^2-x^3+2)(-171x^2-1674x+1000y+1561)
\end{equation}
shows that the map preserves the curve wherever defined. To check
that the required denominators are nonzero as functions, and the
iterates distinct and nonconstant, use the curve point
$P_0=(129/100,-383/1000)$. Eight successive first coordinates are
the following fractions, denoted $r_i$:
\begin{center}
\small
\begin{tabular}{@{}cl@{}}
\toprule
$i$&$r_i=x_i(P_0)$\\
\midrule
0&$\dfrac{129}{100}$\\[3pt]
1&$\dfrac{164323}{29241}$\\[3pt]
2&$\dfrac{2340922881}{58675600}$\\[3pt]
3&$\dfrac{307326105747363}{160280942564521}$\\[3pt]
4&$\dfrac{794845361623184880769}{513127310073606144900}$\\[3pt]
5&$\dfrac{49680317504529227786118937923}{3458519104702616679044719441}$\\[3pt]
6&$\dfrac{30037088724630450803382035538503505921}{3010683982898763071786842993779918400}$\\[3pt]
7&$\dfrac{182386897568483763089689099250940650872600619203}{127572396335305049740547038646345741798859364401}$\\[3pt]
\bottomrule
\end{tabular}
\end{center}
Exact cross multiplication gives
\begin{equation}\label{eq:orbit-order}
 r_0<r_7<r_4<r_3<r_1<r_6<r_5<r_2,
 \qquad r_i\ne3\quad(0\le i\le7).
\end{equation}
The fractions and comparisons are checked by
\ancfile{elliptic.py}. Each iterate is therefore defined as a
rational function, and the first seven $x_i$ are distinct. If
$P_1$ is the image of $P_0$, then $x_i(P_1)=r_{i+1}\ne r_i$
for $0\le i\le6$, so each $x_i$ is nonconstant.

The rational function $x^3-2$ is not a square in $\Q(x)$, by its
odd order at infinity. For every transcendental real
$X>\sqrt[3]{2}$ and either sign $Y=\pm\sqrt{X^3-2}$, substitution
defines an embedding of $F$ into $\R$. Indeed $\Q(x)\simeq\Q(X)$,
and a relation $Y\in\Q(X)$ would make $x^3-2$ a square after
clearing denominators. The quadratic field extension therefore
embeds injectively. Every nonzero denominator and difference of
functions just checked remains nonzero at $(X,Y)$.
The actual circle coordinate field is the embedded copy of $F$:
the coordinates recover $X_0=X$ and $X_1$, and
\[
 Y=\frac{3X^2+9X-4-(X-3)^2X_1}{10}.
\]

\subsection{The tensor ring in this example is a domain}
This property is needed only for the following construction. Put
$K=\Q(x,u)$. The element $x^3-2$ is not a square in $K$.
In $K(y)$, a proposed square root of $u^3-2$ has the form
$a+by$, $a,b\in K$. Comparing the $y$ coefficient in its square
gives $ab=0$. If $b=0$, it would make $u^3-2$ a square in $K$;
if $a=0$, it would make $(u^3-2)/(x^3-2)$ a square in $K$.
Both have odd order at $u=\infty$, so neither is possible.
Consequently $K(y,v)$ with $y^2=x^3-2$, $v^2=u^3-2$
has degree four over $K$.

The ring
\[
 A_0=\Q[x,u,y,v]/(y^2-x^3+2,v^2-u^3+2)
\]
is free over $\Q[x,u]$ with basis $1,y,v,yv$. These elements
are independent in the degree-four field just constructed, so
$A_0$ injects into it and is a domain. The tensor ring
$\cB_F=F\otimes_\Q F$ is the localization of $A_0$ at the
nonzero pure left and pure right elements. Thus it too is a domain.
This argument is specific to this function field; the preceding
classification and two-group certificate assumed no such property.

\subsection{A simultaneous invariant}
For the two copies of the function field write
\[
 Z=(x-u)^2,\qquad S=xu(x+u)-4+2yv,
 \qquad d=S/Z.
\]
The last quotient is taken only in the fraction field of $\cB_F$.
The same map \eqref{eq:elliptic-map} on both pairs $(x,y)$ and
$(u,v)$ preserves $d$. Here is a polynomial verification.
Write
\[
 D=x-3,\quad E=u-3,\quad a=\cA(x,y),\quad a'=\cA(u,v),
 \quad b=\cB(x,y),\quad b'=\cB(u,v),
\]
and put
\begin{equation}\label{eq:invariant-polynomial}
 \begin{split}
 M&=aa'(aE^2+a'D^2)-4D^4E^4+2bb'DE,\\
 \mathcal C&=M(x-u)^2-S(aE^2-a'D^2)^2.
 \end{split}
\end{equation}
Division by the two monic curve equations gives
\begin{equation}\label{eq:invariant-ideal}
 \mathcal C=Q_1(y^2-x^3+2)+Q_2(v^2-u^3+2).
\end{equation}
The quotients are displayed in Section~\ref{sec:certificates}, and
\ancfile{elliptic.py} performs the divisions with exact coefficients.
The transformed $S$ and $Z$ have common denominator $D^4E^4$,
with numerators $M$ and $(aE^2-a'D^2)^2$, respectively.
Thus \eqref{eq:invariant-ideal} proves invariance of $d$.

All the needed differences $L(x_i)-R(x_i)$ are nonzero. In fact
$x_i$ is nonconstant and hence not rational, so there is a
$\Q$-linear functional $\lambda:F\to\Q$ with
$\lambda(1)=0$ and $\lambda(x_i)=1$. Applying
$\lambda\otimes\mathrm{id}$ to a proposed equality
$x_i\otimes1=1\otimes x_i$ gives $1=0$.
The denominators in every $d_i$ are consequently nonzero.
The invariant gives $d_i=d$ along all seven simultaneous iterates.

Set
\[
 K(t,w)=tw(t+w)-4,\qquad C(t,w)=t^2w^2+8(t+w).
\]
The further polynomial identity
\begin{equation}\label{eq:biquadratic}
 K(t,w)^2-(t-w)^2C(t,w)=4(t^3-2)(w^3-2)
\end{equation}
implies, for $t_i=L(x_i)$ and $u_i=R(x_i)$,
\begin{equation}\label{eq:invariant-node}
 (t_i-u_i)^2d^2-2K(t_i,u_i)d+C(t_i,u_i)=0.
\end{equation}
Indeed $S_i=K(t_i,u_i)+2L(y_i)R(y_i)$ and
$d=S_i/(t_i-u_i)^2$; expanding \eqref{eq:biquadratic} proves
\eqref{eq:invariant-node}.

\subsection{A certificate with only pure-factor denominators}
In $\cB_F[t,w]$, define
\begin{equation}\label{eq:elliptic-certificate}
 G(t,w)=-2L((4y^2)^{-2})
 \bigl[(t-w)^2S^2-2K(t,w)ZS+C(t,w)Z^2\bigr].
\end{equation}
This has degree at most two in each variable. Its only inverse
is a pure left factor, since $4y^2\ne0$ in $F$.
Multiplying \eqref{eq:invariant-node} by $Z^2$ proves that
\eqref{eq:elliptic-certificate} vanishes at each tensor node,
first in the fraction field and therefore in the domain $\cB_F$.
On multiplying coefficients, $m(Z)=0$ and $m(S)=4y^2$, so
\[
 m(G(t,w))=-2(t-w)^2.
\]
Equation \eqref{eq:circle-kernel} gives a tensor matrix with
multiplied value $\HH$ and zero node evaluations. The seven points
are distinct, hence span the plane. OpenAI's Theorem~1.1,
sufficient direction, proves the positive Ramsey assertion of
Theorem~\ref{thm:elliptic}.

\subsection{All thirty-five cross ratios are transcendental}
For $0\le a<b<c<d\le6$, put
\[
 \rho_{abcd}=
 \frac{(x_a-x_c)(x_b-x_d)}{(x_a-x_d)(x_b-x_c)}\in F.
\]
At $P_0$ and $P_1$ its values are respectively
$\crat(r_a,r_b;r_c,r_d)$ and
$\crat(r_{a+1},r_{b+1};r_{c+1},r_{d+1})$, where here the
notation denotes the corresponding difference formula.
All differences in the denominators are nonzero by
\eqref{eq:orbit-order}. Computing first value minus second value,
the exact signs in lexicographic order of $(a,b,c,d)$ are
\begin{equation}\label{eq:cross-signs}
 \texttt{--+-++--++-----+--++++---++++---+++}.
\end{equation}
These are $35$ nonzero signs. The rational table, the difference
formula, and \eqref{eq:cross-signs} form a finite certificate;
\ancfile{elliptic.py} enumerates it with \ancfile{fractions.Fraction}
and saves every exact difference.

If $\rho_{abcd}$ were algebraic over $\Q$, its irreducible
polynomial would have a rational root upon evaluation at $P_0$.
This evaluation is legitimate for that finite expression: every
denominator used is nonzero at $P_0$. The irreducible polynomial
would therefore be linear and $\rho_{abcd}$ a rational constant,
contradicting its different value at $P_1$. Thus it is transcendental
in $F$, and injectivity of the real embedding preserves this fact.
All other orderings are invertible fractional linear transforms,
so all four-point cross ratios are transcendental.

Finally $M(A)\subset F$ and $\trdeg_\Q F=1$, while $M(A)$ contains
a transcendental cross ratio. Hence $\trdeg_\Q M(A)=1$.
This completes the proof of Theorem~\ref{thm:elliptic}, with the
positive Ramsey conclusion depending exactly on OpenAI's sufficient
direction as specified.
